\documentclass[10pt]{article}
\usepackage{hyperref}
\usepackage{enumitem}
\usepackage{cmap}
\usepackage[utf8]{inputenc}
\usepackage[T1]{fontenc}
\usepackage{lmodern}
\usepackage[english]{babel}
\usepackage[left=1.05cm,top=0.82cm,right=1.05cm,bottom=0.68cm]{geometry}
\IfFileExists{glyphtounicode.tex}{\input{glyphtounicode}\pdfgentounicode=1}{}
\hypersetup{
    hidelinks,
    pdfauthor={Kevin Bravo},
    pdftitle={Kevin Bravo - Ingeniero de Producto React Next.js TypeScript},
    pdfsubject={CV ATS para un rol de Ingeniero de Producto con React, Next.js, TypeScript y PostgreSQL},
    pdfkeywords={Ingeniero de Producto, Product Engineer, React, Next.js, TypeScript, JavaScript, PostgreSQL, Prisma, Tailwind CSS, SvelteKit, arquitectura frontend, UI responsive, server-side rendering, performance, dashboards, reportes, onboarding, discovery de producto, feedback de usuarios, APIs REST, Django, Node.js, code review, CI/CD, Docker}
}
\setlength{\parindent}{0pt}
\hyphenpenalty=10000
\exhyphenpenalty=10000
\emergencystretch=2em
\pagestyle{empty}
\setlist[itemize]{leftmargin=*, noitemsep, topsep=0pt, partopsep=0pt, parsep=0pt, label={\textbf{-}}}

\newcommand{\sectiontitle}[1]{%
    \vspace{4pt}
    \begin{center}
        \textbf{#1}
    \end{center}
    \vspace{-5pt}
}

\begin{document}
\raggedright

\begin{center}
    {\LARGE \textbf{Kevin Bravo}}\\
    Ingeniero de Producto | React, Next.js y TypeScript \\
    \hrulefill
\end{center}

\begin{center}
    \href{mailto:0bravokevin@gmail.com}{0bravokevin@gmail.com}
    | \href{https://kevinbravo.com}{kevinbravo.com}
    | \href{https://github.com/0bkevin}{github.com/0bkevin}
    | \href{https://www.linkedin.com/in/0bkevin}{linkedin.com/in/0bkevin}
\end{center}

\fontsize{9.45}{10.35}\selectfont

\sectiontitle{Resumen Profesional}
Ingeniero de producto con 5 años de experiencia construyendo y operando productos web full-stack, plataformas internas y dashboards. Trabajo desde los modelos de datos y las APIs hasta la interfaz responsive, el despliegue y el soporte en producción. Lideré la migración de un producto activo de React a SvelteKit en Free2Z y construí plataformas con Next.js y PostgreSQL para AVAA y una unidad del Departamento de Estado de EE.UU. Trabajo cerca de usuarios y stakeholders para hacer que los flujos de producto sean claros, rápidos y útiles en la operación diaria.

\sectiontitle{Competencias Técnicas}
\textbf{Frontend y producto:} React, Next.js, TypeScript, JavaScript, SvelteKit, Tailwind CSS, server-side rendering, UI responsive, manejo de estado, arquitectura frontend, optimización de performance, visualización de datos, discovery de producto y feedback de usuarios. \\
\textbf{Backend y datos:} PostgreSQL, Prisma, SQL, Django, APIs REST, herramientas web de Node.js, modelado relacional, cambios de esquema, migraciones, autenticación, control de acceso por roles, dashboards y sistemas de reportes. \\
\textbf{Calidad y entrega:} Git, code reviews, registros de decisiones técnicas, QA manual, pruebas de regresión, verificación de features, CI/CD, Docker, GCP, Azure, Cloudflare, debugging en producción, ownership de releases y entrega end-to-end.

\sectiontitle{Experiencia Profesional}
\textbf{Proyecto Educa} \hfill \textit{Remoto} \\
\textbf{Chief Technology Officer} \hfill \textit{Dic 2025 - Presente}
\begin{itemize}
    \item Defino la dirección de producto y técnica de plataformas educativas, llevando las necesidades de stakeholders desde la arquitectura y las prioridades de implementación hasta los releases y el soporte en producción.
    \item Reviso el trabajo de contribuidores y las decisiones técnicas, establezco estándares de release y desbloqueo la implementación entre frontend, backend y entrega.
\end{itemize}

\vspace{2pt}
\textbf{Free2Z} \hfill \textit{Remoto} \\
\textbf{Principal Software Engineer} \hfill \textit{Abr 2025 - Jun 2026}
\begin{itemize}
    \item Lideré la migración de un producto activo de React a SvelteKit, mejorando la arquitectura frontend, SSR, manejo de estado, mantenibilidad y performance para un producto con usuarios reales.
    \item Construí features conectadas al backend con APIs de Python y Django, video en vivo con Dyte, Tailwind CSS, un editor Markdown, un CMS propio, publicación, búsqueda y flujos operativos.
    \item Integré y publiqué un asistente de escritura con IA en el flujo de blogging, ayudando a los usuarios a redactar posts y elegir temas a partir de la actividad de mayor interés en la plataforma.
\end{itemize}

\vspace{2pt}
\textbf{Departamento de Estado de EE.UU. (vía AVAA)} \hfill \textit{Remoto} \\
\textbf{Product Engineer | Billingua Talent} \hfill \textit{Nov 2024 - Oct 2025}
\begin{itemize}
    \item Construí Billingua Talent como único ingeniero, siendo responsable del modelo de datos relacional, los flujos sobre PostgreSQL, el acceso por roles, las APIs, dashboards, reportes, despliegue y mantenimiento.
    \item Realicé demos y un piloto con 50 usuarios y 5 empresas; después ajusté el registro, la gestión de servicios, la validación de datos y los reportes a partir del feedback de usuarios.
\end{itemize}

\vspace{2pt}
\textbf{Freelance} \hfill \textit{Remoto} \\
\textbf{Software Engineer} \hfill \textit{Sep 2024 - Dic 2024}
\begin{itemize}
    \item Entregué productos web para clientes, proyectos frontend, scripts de automatización y flujos conectados a APIs, ocupándome del discovery, la implementación, el despliegue, la iteración y el soporte.
\end{itemize}

\vspace{2pt}
\textbf{Asociación Venezolano Americana de Amistad (AVAA)} \hfill \textit{Remoto} \\
\textbf{Software Developer} \hfill \textit{Sep 2022 - Jul 2024} \\
\textbf{Software Developer Intern} \hfill \textit{Sep 2021 - Ago 2022}
\begin{itemize}
    \item Construí, desplegué y mantuve SEP con Next.js, PostgreSQL, Prisma, Auth.js y Azure, reemplazando hojas de cálculo por una fuente de verdad multi-tenant para aproximadamente 300 becarios en 3 capítulos.
    \item Construí SEA para EducationUSA, centralizando historiales de estudiantes, actividades de seguimiento, servicios, pagos, recibos, relaciones institucionales y reportes en una plataforma full-stack.
\end{itemize}

\sectiontitle{Proyectos Seleccionados de Producto y Herramientas}
\textbf{Zrode y Brio} \hfill \textit{Proyectos independientes} \\
Construyo y uso un entorno de desarrollo agéntico y un control plane móvil para flujos de desarrollo, pairing, control de sesiones, relay routing, logs de ejecución, health checks y visibilidad operativa. \\
\textbf{Código abierto y liderazgo:} Contribuí al editor de código open source Zed; soy Vicepresidente de JA Venezuela Alumni; organicé un hub para un hackathon global de Web3 e IA; y acompañé a estudiantes desde la investigación de producto hasta un prototipo móvil funcional.

\sectiontitle{Educación e Idiomas}
\textbf{Instituto de Estudios Superiores de Administración (IESA)} \hfill \textit{2025 - 2027} \\
Maestría en Administración Pública, en curso. \textbf{Idiomas:} Español (Nativo), Inglés (competencia profesional C1).

\end{document}
